%%%%%%%%%%%%%%%%%%%%%%%%%%%%%%%%%
%
%       EXERCÍCIO
%
%%%%%%%%%%%%%%%%%%%%%%%%%%%%%%%%%

\ifdefstring{\atividade}{prova}{%
    \ifdefstring{\modo}{objetivo}{%
        \renewcommand{\valorquestao}{\ValorQObj\ ponto}
    }{%
        \renewcommand{\valorquestao}{\ValorQDisc\ pontos}
    }%
}%

\begin{exercicioBanco}[\valorquestao]
Quais das funções a seguir são função quadrática?
\begin{multicols}{3}
    \begin{enumerate}[(I)]
        \item \(f(x) = 2 - x\).
        %
        \item \(f(x) = (x + 1)^{2}\).
        %
        \item \(f(x) = |x^{2} - 1|\).
        %
        \item \(f(x) = x(x - 1)\).
        %
        \item \(f(x) = 2 - x + x^{3}\).
        %
        \item \(f(x) = \pi x^2\).
    \end{enumerate}
\end{multicols}

% Define as alternativas
\newcommand{\alternativas}{%
    \begin{center}
        \begin{tabularx}{\textwidth}{XXXXX}
            (a) I, II e IV. &
            (b) II e IV. &
            (c) II, IV e VI. &
            (d) I, II, III, IV e VI. &
            (e) II, III e VI.
        \end{tabularx}
    \end{center}
}

% Define a resposta correta
\newcommand{\resposta}{C}

% Lógica condicional para exibição
\ifdefstring{\atividade}{lista}{%
    \alternativas
    \vspace{0.5em}
    
    \noindent\textbf{Resposta:} letra \textbf{\resposta}.
}{%
    \ifdefstring{\modo}{objetivo}{%
        \alternativas
    }{%
        % modo = discursiva → não mostra alternativas
    }
}
\end{exercicioBanco}

